\begin{onehalfspace}

Il Capitolo~\ref{chap:risultati} ha chiuso la verifica dell'ipotesi iniziale con un esito netto: applicare l'apprendimento automatico \emph{a valle} dell'esecuzione, sul solo output, non ha senso --- l'informazione contenuta nell'istogramma è già interamente sfruttabile da una ricerca multi-candidato a costo zero. Questo capitolo presenta il nucleo della tesi, che nasce proprio da quel risultato: se l'\ac{ML} deve contribuire alla mitigazione degli errori dell'algoritmo di Shor, deve operare \emph{a monte}, sull'informazione che nell'output non c'è --- la fisica del dispositivo. Nelle sezioni che seguono si descrive il problema (gli errori hardware non sono uniformi, ma sistematici e quindi prevedibili), la soluzione oggi adottata dall'industria (la duplicazione dei qubit), la proposta di questa tesi (un modello predittivo degli errori addestrato sulla specifica \mbox{QPU}) e l'architettura del sistema con cui intendiamo realizzarla e validarla.

% -----------------------------------------------
\section{Il Problema: gli Errori Hardware non sono Uniformi}
\label{sec:errori_non_uniformi}

Il Capitolo~\ref{chap:rumore} ha analizzato le quattro fonti fisiche di errore dei processori \ac{NISQ}: la decoerenza ($T_1$, $T_2$), gli errori di gate, gli errori di misura (\textit{readout}) e il crosstalk. Nella modellazione adottata finora --- come nella maggior parte della letteratura introduttiva --- questi errori sono stati trattati come proprietà \emph{globali} del dispositivo: un unico $T_1$, un unico tasso $\varepsilon_{2q}$, un'unica probabilità di errore di lettura per tutti i qubit. La realtà fisica di una \ac{QPU} è diversa, e questa differenza è il punto di partenza dell'intero approccio proposto.

\paragraph{La decoerenza varia da qubit a qubit.}
I qubit superconduttori sono strutture litografiche a scala molecolare: imperfezioni nel processo di fabbricazione (ad esempio nell'ossidazione delle giunzioni Josephson) fanno sì che, sullo stesso chip, un qubit possa esibire $T_1 = 100\,\mu\text{s}$ e il suo vicino $T_1 = 30\,\mu\text{s}$. La conseguenza operativa è severa: quando due qubit vengono posti in entanglement, la coerenza della coppia decade con il \emph{peggiore} dei due --- il qubit da $100\,\mu\text{s}$ viene trascinato a $30\,\mu\text{s}$, e quando uno dei due decade, l'informazione condivisa è persa per entrambi. Un algoritmo che ignora questa disomogeneità è costretto a dimensionarsi sul qubit più debole della linea utilizzata; il rischio cresce quando il tempo di esecuzione del circuito si avvicina al bordo della finestra di coerenza, dove la probabilità di decadimento è massima.

\paragraph{I gate sono hardware, e ogni linea è diversa.}
Le porte logiche quantistiche non risiedono nel chip: sono impulsi a microonde generati dall'elettronica di controllo esterna e convogliati ai qubit tramite guide d'onda. Lo stesso gate logico, applicato a linee diverse, può richiedere tempi diversi ed esibire fedeltà diverse; un impulso impreciso ruota lo spin di un angolo sbagliato, producendo uno stato errato che si manifesta solo alla misura. Il tempo di gate --- decine di nanosecondi sui superconduttori, millisecondi sugli ioni intrappolati --- è il secondo parametro fondamentale insieme alla coerenza: il numero massimo di operazioni eseguibili è il rapporto tra la finestra di coerenza e la durata del singolo gate, ed è questo rapporto (non il solo $T_1$) a determinare se una \ac{QPU} può fisicamente sostenere il circuito di Shor per una data istanza.

\paragraph{La misura è essa stessa una sorgente di errore.}
La lettura di un qubit avviene bombardandolo con una microonda e osservando la risposta: è un processo fisico che distrugge lo stato e che, se il qubit è prossimo alla fine della coerenza, può restituire 0 per 1 o viceversa. Questi errori si caratterizzano tramite matrici di confusione ottenute preparando stati noti e misurandone le distribuzioni di lettura. Nei dispositivi commerciali l'affidabilità di lettura considerata accettabile parte dal 95\%: un errore intrinseco superiore al 5\% rende il dato inutilizzabile.

\paragraph{Il crosstalk dipende dalla topologia.}
I qubit adiacenti sul chip si accoppiano capacitivamente: operare su un qubit perturba i vicini con rotazioni parassite non intenzionali, e la stessa microonda di controllo o di lettura, per quanto focalizzata, raggiunge parzialmente i qubit circostanti. L'entità del fenomeno dipende dalla \emph{geometria} del dispositivo --- una catena lineare, una griglia, una connettività completa --- e dalle linee di accoppiamento fisico effettivamente presenti: su una topologia a catena, l'entanglement tra qubit non adiacenti richiede operazioni intermedie che amplificano il disturbo e accorciano la vita dello stato condiviso.

\paragraph{Sistematicità, quindi prevedibilità.}
Il punto cruciale è che questi difetti non sono rumore bianco: sono \textbf{caratteristiche stabili del singolo esemplare fisico}. Se una linea di qubit fallisce otto volte su dieci per circuiti di una certa profondità, quella linea è \emph{sistematicamente} debole per quella classe di input --- e continuerà a esserlo. La profondità del circuito di Shor dipende dall'istanza ($N$ piccolo usa meno gate di $N$ grande), per cui una stessa linea può essere adeguata per certi input e inadeguata per altri. Un comportamento sistematico e dipendente da variabili osservabili (quale linea, quale input, quale profondità) è, per definizione, un comportamento \emph{apprendibile}.

% -----------------------------------------------
\section{La Pratica Attuale: la Duplicazione dei Qubit}
\label{sec:duplicazione}

Come gestisce oggi l'industria l'inaffidabilità dei singoli qubit? Il teorema di no-cloning (Capitolo~\ref{chap:fondamenti}) impedisce la soluzione classica più ovvia: non è possibile copiare uno stato quantistico durante l'esecuzione, né "trasferire" il calcolo da un qubit morente a uno sano. La tecnica adottata è la \textbf{duplicazione}: per ogni qubit logico dell'algoritmo se ne impiegano due fisici, e l'intero calcolo viene raddoppiato.

Operativamente, si eseguono \emph{due istanze gemelle} dell'algoritmo di Shor in parallelo, sugli stessi dati, su due linee di qubit diverse del dispositivo: i registri iniziali vengono preparati in modo identico ($A_1$/$A_2$, $B_1$/$B_2$, \ldots), le due istanze evolvono indipendentemente, e alla fine si leggono gli stati di entrambe. È più probabile che, su due copie, almeno una arrivi viva alla misura: dove una copia risulta decaduta o illeggibile, si utilizza l'informazione dell'altra. Poiché le due istanze operano sugli stessi input, non possono produrre soluzioni valide diverse: se entrambe sopravvivono concordano, se una muore si usa la superstite, se muoiono entrambe l'esecuzione è da ripetere.

Il costo di questa strategia è diretto e severo: \textbf{il fabbisogno di qubit fisici raddoppia}. Per l'istanza crittograficamente rilevante --- fattorizzare un modulo RSA a 2048 bit, che richiede $\sim$2048 qubit logici --- servono $\sim$4096 qubit fisici. L'esperienza pratica mostra inoltre rendimenti decrescenti: raddoppiare conviene, triplicare o quadruplicare no --- il guadagno di affidabilità aggiuntivo non giustifica il costo, proprio perché il meccanismo agisce solo sulla probabilità di sopravvivenza e non sulla \emph{causa} dei fallimenti.

La duplicazione ha infine tre limiti strutturali che la rendono un bersaglio naturale per un approccio più intelligente:
\begin{itemize}
    \item è \textbf{agnostica rispetto al dispositivo}: tratta tutti i qubit come ugualmente inaffidabili, ignorando che i difetti sono localizzati e stabili;
    \item \textbf{spreca conoscenza}: ogni esecuzione rivela quali linee hanno fallito, ma questa informazione non viene accumulata né riutilizzata;
    \item \textbf{consuma metà della QPU}: la capacità computazionale effettiva del dispositivo è dimezzata in modo permanente, indipendentemente da quanto la specifica esecuzione ne avrebbe avuto bisogno.
\end{itemize}

% -----------------------------------------------
\section{La Proposta: un Modello Predittivo degli Errori}
\label{sec:proposta_predittiva}

La proposta di questa tesi è sostituire la ridondanza fisica con la \textbf{conoscenza appresa del dispositivo}: un modello di apprendimento automatico, addestrato sul comportamento della specifica \ac{QPU}, che predica dove e quando gli errori si manifesteranno.

\paragraph{Perché qui l'ML ha senso (e a valle no).}
La verifica del Capitolo~\ref{chap:risultati} ha dimostrato che un classificatore che osserva il solo istogramma di output non aggiunge nulla: quell'informazione è già catturata dalla ricerca TOP-K. Il modello proposto attinge invece a informazione \emph{strutturale}, che nessuna strategia a valle può vedere: la topologia del dispositivo, i tempi di coerenza dei singoli qubit, quali linee il circuito attraversa, la profondità richiesta dall'input. È la differenza tra giudicare una fotografia sfocata e conoscere quale lente dell'obiettivo è graffiata: la seconda conoscenza permette di \emph{prevedere} la sfocatura prima dello scatto, e di scegliere una lente diversa.

\paragraph{Che cosa promette il modello.}
Il livello minimo di funzionalità --- sufficiente a giustificare l'approccio --- è la \textbf{segnalazione}: dato l'esito di un'esecuzione e le condizioni in cui è avvenuta, il modello indica "questo risultato è probabilmente errato". Non serve conoscere la risposta giusta: la sola segnalazione affidabile permette di scartare l'esito e reiterare, esattamente come la strategia TOP-K permetteva di evitare iterazioni sprecate. Il livello avanzato è la \textbf{prescrizione}: per un dato tipo di input, il modello indica quali linee di qubit evitare ("per questa profondità di circuito, la linea 1 fallisce sistematicamente: usa la seconda"), instradando l'esecuzione sulle risorse fisiche più adatte \emph{prima} di consumarle.

\paragraph{Un modello per esemplare, non per modello.}
Poiché i difetti sono caratteristiche del singolo chip --- due esemplari identici dello stesso modello di QPU hanno mappe di difetti diverse --- il modello va addestrato \emph{su quella specifica QPU}, con una fase di caratterizzazione basata su input a verità nota. Questo è un costo una-tantum per dispositivo, analogo alla calibrazione periodica che i provider già eseguono, ma il suo prodotto è riutilizzabile per ogni esecuzione successiva.

\paragraph{Il claim da validare.}
La promessa quantitativa, che la fase sperimentale dovrà verificare, si formula così: \emph{l'algoritmo di Shor eseguito in copia singola, assistito dal modello predittivo, raggiunge un'affidabilità almeno pari a quella della duplicazione, utilizzando la metà dei qubit fisici}. Se il claim regge, il messaggio pratico è quello indicato dal relatore: chi implementa Shor non ha più bisogno di duplicare --- esegue una copia sola e lascia che il modello, addestrato su quella macchina, faccia il resto.

% -----------------------------------------------
\section{Architettura del Sistema Proposto}
\label{sec:architettura_predittiva}

Il sistema si compone di quattro elementi, progettati per essere realizzati con l'infrastruttura sperimentale già sviluppata (Capitoli~\ref{chap:strumenti} e~\ref{chap:sviluppo}) opportunamente estesa.

\paragraph{1. QPU virtuale con iniezione di errori configurabile.}
Il primo componente è un simulatore di \ac{QPU} in cui le quattro categorie di errore sono configurabili \emph{per singolo qubit e per singola linea}, non globalmente: a ogni qubit i propri $T_1$/$T_2$, a ogni coppia accoppiata la propria fedeltà e durata di gate, a ogni linea di lettura il proprio errore di misura, e un modello di crosstalk derivato dalla topologia dichiarata (quali qubit sono fisicamente accoppiati con quali). Rispetto al noise model uniforme utilizzato nella fase di verifica, questa è un'estensione sostanziale ma tecnicamente percorribile: Qiskit Aer supporta nativamente errori per-qubit e coupling map, mentre il crosstalk andrà modellato come canale di errore correlato sui vicini topologici. La QPU virtuale si comporta come una QPU reale --- "su dieci gate che passano, a volte uno sbaglia" --- ma con percentuali e localizzazione dei difetti decise da noi, e quindi note.

\paragraph{2. Dataset a verità nota.}
Il secondo componente è la campagna di caratterizzazione: si costruisce un insieme di istanze di fattorizzazione con soluzione nota (semiprimi $N = p \times q$ con $p$, $q$ primi scelti, di profondità circuitale crescente) e lo si esegue ripetutamente sulla QPU virtuale --- nell'ordine del migliaio di esecuzioni per configurazione. Poiché input, output atteso e profilo di errore iniettato sono tutti noti, ogni esecuzione produce un esempio etichettato: quali linee erano coinvolte, quale profilo di errore era attivo, e se l'esito era corretto.

\paragraph{3. Modello predittivo con feature strutturali.}
Il terzo componente è il modello di apprendimento: le feature comprendono la descrizione della QPU (topologia, $T_1$/$T_2$ per qubit, fedeltà di gate per linea), la descrizione dell'esecuzione (linee utilizzate, profondità del circuito, tipo di input) e l'esito osservato; le etichette derivano dalla verità nota. L'architettura del modello (classificatore, regressore di affidabilità per linea, o sistema a due stadi) è una scelta progettuale aperta, da guidare con lo stesso rigore comparativo usato nella fase di verifica. Un'estensione avanzata segnalata dal relatore --- l'uso delle matrici di confusione come spazio hamiltoniano in ingresso a un algoritmo \ac{QAOA} per identificare gli errori trascurabili --- è deliberatamente esclusa dal perimetro di questa fase per la sua complessità, e rimandata agli sviluppi futuri.

\paragraph{4. Protocollo di validazione con baseline di duplicazione.}
Il quarto componente incorpora la lezione metodologica della fase di verifica: \emph{l'ablazione si progetta dall'inizio, non si aggiunge dopo}. Il confronto sperimentale prevede tre bracci fin dal disegno: (a) copia singola senza modello (baseline inferiore), (b) \textbf{duplicazione dei qubit} --- da implementare come baseline di riferimento, poiché è la pratica che il modello intende sostituire --- e (c) copia singola assistita dal modello predittivo. Le metriche sono il tasso di successo, il numero di iterazioni e il costo in qubit fisici; la significatività statistica sarà valutata con gli stessi strumenti non parametrici già adottati (Mann-Whitney, $K \geq 30$ ripetizioni).

% -----------------------------------------------
\section{Ipotesi Operative e Criteri di Successo}
\label{sec:ipotesi_predittiva}

La fase sperimentale è guidata da due ipotesi verificabili:

\begin{enumerate}
    \item \textbf{Ipotesi di prevedibilità}: sugli errori generati dalla QPU virtuale con difetti localizzati e stabili, un modello addestrato con feature strutturali raggiunge una capacità di segnalazione degli esiti errati significativamente superiore al caso ($\text{AUC} > 0.5$ con margine statisticamente significativo), e crescente con la sistematicità dei difetti iniettati.
    \item \textbf{Ipotesi di sostituzione}: la copia singola assistita dal modello eguaglia o supera il tasso di successo della duplicazione a parità di iterazioni, con la metà dei qubit fisici. Criterio quantitativo: differenza di tasso di successo non inferiore (test unilaterale, $p < 0.05$) rispetto al braccio duplicato.
\end{enumerate}

Diversamente dalla fase di verifica --- dove le metriche di classificazione eccellenti nascondevano l'assenza di guadagno operativo --- il criterio di successo primario è qui direttamente operativo (ipotesi 2), e le metriche di apprendimento (ipotesi 1) hanno il solo ruolo di diagnostica intermedia. La realizzazione sperimentale di questo programma --- l'implementazione della QPU virtuale, la generazione del dataset, l'addestramento e la validazione comparativa --- costituisce la prossima fase del lavoro; le estensioni oltre questo perimetro, dalla validazione su hardware fisico reale all'integrazione con la correzione d'errore strutturale, sono discusse nel Capitolo~\ref{chap:sviluppi}.

\end{onehalfspace}
